\clearpage
\fancyhead[L]{\small DELAMBRE. TABLES ÉCLIPTIQUES.}
\hypertarget{native-608-lower}{}
\begin{center}\small Göttingische gelehrte Anzeigen. Stück 45. Seite 441..443. 1818 März 19.\\[3pt]\rule{.72\linewidth}{.3pt}\end{center}

\emph{Tables écliptiques des satellites de jupiter, d'après la théorie de M. le marquis de \textup{\textsc{Laplace}}, et la totalité des observations faites depuis \textup{1662} jusqu'à l'an \textup{1802}; par M. \textup{\textsc{Delambre}}, chevalier de St. Michel etc.} Paris 1817. Bei der Wittwe Courcier. Die Einleitung 58 Seiten, die Tafeln 32 halbe Bogen. 4.

Die vor beinahe 30 Jahren vom Hrn. \textsc{Delambre} nach \textsc{Laplace}'s Theorie berechneten, und in der dritten Ausgabe von \textsc{Lalande}'s Astronomie abgedruckten Tafeln für die Verfinsterungen der Jupiterstrabanten, sind bekanntlich von den Astronomen mit verdientem Beifall aufgenommen und in allgemeinen Gebrauch gekommen. Jener arbeitsame und unverdrossene Rechner hat seitdem die vervielfältigten Beobachtungen und die noch grössere Vervollkommnung der Theorie benutzt, um die Grundlagen seiner Tafeln von neuem zu verbessern, und so sind die vorliegenden neuen Tafeln entstanden, welche bereits vor mehrern Jahren gedruckt, aber erst jetzt ausgegeben sind. Der grosse Geometer, unter dessen Leitung diese mühsame Arbeit ausgeführt ist, hat ihr bereits das ehrenvolle Zeugniss gegeben: `\textsc{Delambre} a exécuté ce travail important avec le plus grand succès; et ses tables qui réprésentent les observations avec l'exactitude des observations mêmes offrent au navigateur un moyen sûr et facile pour avoir sur le champ, par les eclipses des satellites, et surtout par celles du premier, la longitude des lieux où il atterre.' Die Einrichtung der Tafeln ist wenig von der der frühern verschieden, und die Einleitung gibt eine sehr ausführliche Anleitung zum Gebrauche derselben. Gewünscht hätten wir nur, dass der Verf. von den Resultaten seiner Verbesserungsarbeit mehr Detail mitgetheilt hätte, so dass man in den Stand gesetzt wäre, die Veränderungen, welche die einzelnen constanten durch die Beobachtungen auszumittelnden Grundlagen (deren Anzahl, die Geschwindigkeit des Lichts mitgezählt, sich be\SourceLeaf{609}kanntlich auf zwei und dreissig beläuft) erlitten haben, klar und leicht zu übersehen, und den Grad
der ihnen beizulegenden Genauigkeit zu beurtheilen. Aus einer Stelle der Einleitung lässt sich schliessen,
dass diese Veränderungen alle sehr gering gewesen sind. Er erklärt, dass obgleich die Fehler der neuen
Tafeln im Allgemeinen geringer sind, als die der frühern `ce que nous avons gagné ne vaut peut-être
pas le travail qu'il a coûté', ganz im Geiste der Ansicht, den Werth einer solchen Arbeit nur nach
dem practischen Nutzen für die geographischen Längenbestimmungen zu würdigen. Nur von dem Coëfficienten
der Lichtgleichung, welchen Hr. \textsc{Delambre} zu $493''2$ annimmt, bemerkt er, dass diese Bestimmung
sich auf mehr als 1000 Beobachtungen des ersten Trabanten gründe, und dass die nicht merklich
davon verschiedene frühere Bestimmung auf 500 Beobachtungen beruhet habe. Es ist doch sehr merkwürdig,
dass dieser Coëfficient die Aberrationsconstante $20''25$ gibt, während dieselbe aus den Rectascensionen
des Polarsterns bestimmt, etwas grösser ausfällt; es möchte indessen zu voreilig sein, hieraus
schon auf eine Verschiedenheit der Geschwindigkeit des Fixsternlichts und des Trabantenlichts, oder
auf eine verschiedene Geschwindigkeit des Lichts am Umfang der Erdbahn von der näher nach der
Sonne schliessen zu wollen. Doch verdient dieser Umstand gewiss fortgesetzte Aufmerksamkeit.
